% ---------------------------------------------------------------------
% MỤC LỤC, DANH SÁCH HÌNH, BẢNG, KÝ HIỆU
% ---------------------------------------------------------------------
\renewcommand{\contentsname}{Mục lục}
\phantomsection\addcontentsline{toc}{chapter}{MỤC LỤC}
\markboth{MỤC LỤC}{}
\tableofcontents

\renewcommand{\listfigurename}{Danh sách hình vẽ}
\phantomsection\addcontentsline{toc}{chapter}{DANH SÁCH HÌNH VẼ}
\listoffigures

\renewcommand{\listtablename}{Danh sách bảng}
\phantomsection\addcontentsline{toc}{chapter}{DANH SÁCH BẢNG}
\listoftables

\unnumchapter{DANH SÁCH CÁC KÝ HIỆU, CHỮ VIẾT TẮT}
\begin{longtable}{L{2.6cm}L{11.5cm}}
\toprule
\textbf{Ký hiệu} & \textbf{Ý nghĩa}\\
\midrule\endhead
AI      & Artificial Intelligence -- Trí tuệ nhân tạo\\
API     & Application Programming Interface -- Giao diện lập trình ứng dụng\\
BN      & Batch Normalization -- Chuẩn hóa theo lô\\
CI/CD   & Continuous Integration/Continuous Deployment -- Tích hợp và triển khai liên tục\\
CWD     & Channel-Wise Distillation -- Chưng cất tri thức theo kênh\\
EC2     & Amazon Elastic Compute Cloud -- Dịch vụ máy ảo của AWS\\
GHCR    & GitHub Container Registry -- Kho image container của GitHub\\
KD      & Knowledge Distillation -- Chưng cất tri thức\\
PPE     & Personal Protective Equipment -- Trang bị bảo hộ lao động\\
Q/DQ    & Quantize/Dequantize -- Cặp nút lượng tử hóa/giải lượng tử trong ONNX\\
SHO     & Sea-Horse Optimizer -- Thuật toán tối ưu cá ngựa\\
SSE     & Server-Sent Events -- Cơ chế đẩy sự kiện từ máy chủ qua HTTP\\
WS      & WebSocket -- Giao thức truyền hai chiều trên một kết nối TCP\\
CUDA    & Compute Unified Device Architecture -- Nền tảng tính toán GPU của NVIDIA\\
CPU     & Central Processing Unit -- Bộ xử lý trung tâm\\
DDD     & Domain-Driven Design -- Thiết kế hướng miền\\
GPU     & Graphics Processing Unit -- Bộ xử lý đồ họa\\
NPU     & Neural Processing Unit -- Bộ xử lý mạng nơ-ron\\
Edge AI & Thực thi tác vụ trí tuệ nhân tạo trên hoặc gần nguồn sinh dữ liệu\\
FP16    & Floating Point 16-bit -- Số thực dấu phẩy động 16 bit\\
FP32    & Floating Point 32-bit -- Số thực dấu phẩy động 32 bit\\
FPS     & Frames Per Second -- Số khung hình xử lý trong một giây\\
INT8    & Integer 8-bit -- Số nguyên 8 bit\\
IoU     & Intersection over Union -- Tỷ lệ giao trên hợp\\
LLM     & Large Language Model -- Mô hình ngôn ngữ lớn\\
mAP     & mean Average Precision -- Độ chính xác trung bình trên các lớp\\
MCP     & Model Context Protocol -- Giao thức kết nối mô hình với công cụ và dữ liệu\\
MQTT    & Message Queuing Telemetry Transport -- Giao thức truyền thông điệp nhẹ\\
ONNX    & Open Neural Network Exchange -- Định dạng trao đổi mô hình học máy\\
PTQ     & Post-Training Quantization -- Lượng tử hóa sau huấn luyện\\
QAT     & Quantization-Aware Training -- Huấn luyện có nhận biết lượng tử hóa\\
RAG     & Retrieval-Augmented Generation -- Sinh nội dung có tăng cường truy xuất\\
RAM     & Random Access Memory -- Bộ nhớ truy cập ngẫu nhiên\\
RBAC    & Role-Based Access Control -- Kiểm soát truy cập theo vai trò\\
REST    & Representational State Transfer -- Phong cách kiến trúc dịch vụ web\\
TLS     & Transport Layer Security -- Giao thức bảo mật tầng truyền tải\\
YOLO    & You Only Look Once -- Họ mô hình phát hiện đối tượng một giai đoạn\\
\bottomrule
\end{longtable}
